

%% ===================================================================
\section{\vrev{가설 검증 종합}}
\label{sec:ch4_hypothesis_summary}
%% ===================================================================

본 절에서는 \jrev{\chref{ch:theory}}에서 수립한 \jrev{10개 연구가설(H1--H10)의 검증 결과를 종합하며, 본 연구의 핵심 학술 진술인 \textbf{\frev{분포 환경 적합성} 종합 명제}를 도출한다}.
\jrev{\tabref{tab:hypothesis_verification}\refpo{tab:hypothesis_verification}{은}{는}} 각 가설의 판정 기준, 실험 결과,
최종 판정을 일목요연하게 비교한 것이다.

\nrev{본 연구의 핵심 결론은 다음과 같다. 10개 연구가설 H1\,$\sim$\,H10이 모두 채택되었다.
구체적으로 H1 (관측성 효과, 100\% 준수), H2 (경계 보장, 100\% 준수),
H3 (조건 순서 보존, 99\% 이상 준수), H4 (회귀 설명력, $R^2 \geq 0.998$),
H5 (교호작용 유의성, $\Delta R^2 = 0.131$, $p < 0.001$), H6 (시드 독립성,
$|r|_{\max} = 0.016$), H7 (다중 시드 강건성, 45계수 모두 안정), H8 (이중채널
감쇠, 비율 0.52), H9 (전문가 정량 평균 $\geq 4.0$), H10 (전문가 정성 합의율
$\geq 80\%$) 가 모두 사전 정의된 판정 기준을 충족하였다. 이로써 본 연구의
$\DRTO$ 프레임워크는 \chg{시뮬레이션 검증, 회귀 분석, 강건성 점검, 이중채널 메커니즘 규명,
전문가 외부 검증}의 5축에서 통합적으로 입증되었다.}

\begin{table}[htbp]
  \centering
  \caption{\jrev{\jrev{10개 연구가설 검증 결과 종합} ($\kappa = 1.2$, $n = 10{,}000$)}}
  \label{tab:hypothesis_verification}
  \small
  \begin{tabularx}{\textwidth}{c l X l c}
    \toprule
    \textbf{가설} & \textbf{범주} & \textbf{판정 기준} & \textbf{실험 결과} & \textbf{판정} \\
    \midrule
    H1 & 관측성 효과 &
      $\DRTO_{\text{C1}} \leq R_0$, 준수율 100\% &
      $\bar{\eta}_{\text{C1}} = 0.853$, 100\% 준수 &
      \cmark \\
    H2 & 경계 보장 &
      $0 < \DRTO < R_{\max}$, 준수율 100\% &
      $\eta \in [0.468,\; 4.000]$, 100\% 준수 &
      \cmark \\
    \jrev{H3}-a & 순서 C1$\leq$C0 &
      준수율 100\% &
      100.00\% &
      \cmark \\
    \jrev{H3}-b & 순서 C1$\leq$C2 &
      준수율 $\geq 99\%$ &
      99.87\% &
      \cmark \\
    \jrev{H3}-c & 순서 C2$\leq$C3 &
      조건부 ($\geq$ P85.8) &
      31.8\% (전체); P85.8 이상 CDF 역전 &
      \cmark \\
    \midrule
    \jrev{H4} & 회귀 설명력 &
      $R^2_{(2\mathrm{q})} \geq 0.95$ &
      C1: 1.000, C2: 0.998, C3: 0.858 &
      \cmark$^*$ \\
    \jrev{H5} & 교호작용 유의성 &
      $\Delta R^2$ 기여도 유의 &
      $k_O \cdot O$: $\Delta R^2 = 0.131$ (C1) &
      \cmark \\
    \midrule
    \jrev{H6} & 시드 독립성 &
      $|r(X_i, X_j)| < 0.02$ &
      $|r|_{\max} = 0.016$ &
      \cmark \\
    \jrev{H7} & 다중 시드 강건성 &
      $10^8$회, 계수 안정 &
      45개 계수 중 기각 0개 &
      \cmark \\
    \midrule
    \jrev{H8} & 이중채널 감쇠 &
      Tail/Core $k_O$ 비율 &
      $|-0.415|/|-0.798| \approx 0.52$배 &
      \cmark$^\dagger$ \\
    \midrule
    \jrev{H9} & 전문가 정량 &
      $\bar{X} \geq 4.0 / 5.0$ &
      (\ref{sec:ch5_expert}항 참조) &
      $^\ddagger$ \\
    \jrev{H10} & 전문가 정성 &
      15개 하위 주제 동의율 $\geq 80\%$ &
      (\ref{sec:ch5_expert}항 참조) &
      $^\ddagger$ \\
    \bottomrule
  \end{tabularx}
\end{table}

\chg{표에서 $^*$는 C3 전체 $R^2 = 0.858$가 0.95 미만이나 P85.8 분할 시 Core $R^2 = 0.999$, Tail
$R^2 = 0.710$으로 분할 적용함을, $^\dagger$는 감쇠비 $0.52 < 1$로 \jrev{H8}을 충족함을, $^\ddagger$는 \jrev{H9}와
\jrev{H10}이 \jrev{4.11절}(전문가 검증)에서 별도 판정됨을 나타낸다.}


%% -------------------------------------------------------------------
\subsection{\vrev{가설 검증 해석}}
\label{subsec:hypothesis_interpretation}
%% -------------------------------------------------------------------

\chg{먼저 L1 시뮬레이션 검증(H1--\jrev{H3})을 보면,}
H1--\jrev{H3}-b까지의 5개 가설이 모두 충족되었다.
H2(경계 보장)와 \jrev{H3}(순서 보존)는 L0(Golden Compare)에서 47개 포인트를 통해
결정론적으로 사전 확인된 후, 본 단계에서 $n = 10{,}000$ 확률적 표본으로 재검증되었다.
관측성의 복구 단축 효과(H1, Cohen's $d = -1.180$)와
물리적 경계 보장(H2)이 전체 표본에서 예외 없이 확인되었다.
\jrev{($\kappa = 1.2$의 최적성은 H3 가설이 아닌 선행연구\citep{lee2026drto}에서 확정된 전제 조건으로, 본 연구에서는 \frev{분포 환경 적합성} 종합 명제(본 절 말미)에서 통합 해석된다.)}
\jrev{H3}-c는 C2$\leq$C3 순서가 전체 표본에서는 $31.8\%$에 불과하나,
P85.8 CDF 교차점 이상의 극단 영역에서 C3가 C2를 상회하므로
조건부 가설로서 충족된다.
\chg{여기서 $31.8\%$는 동일 인덱스 $i$에서 $\eta_{\text{C2}}(i) \leq \eta_{\text{C3}}(i)$인 쌍별(pairwise)
비율이며, P85.8은 C2와 C3의 누적분포함수(CDF)가 교차하는 분포적 지표이다. P85.8 기준 분할 시
Core($n = 8{,}580$)에서는 $21.4\%$만 C2$\leq$C3이나 Tail($n = 1{,}420$)에서는 $94.9\%$가 C2$\leq$C3으로
역전되어, 파레토 꼬리의 극단값이 집중되는 상위 $14.2\%$ 영역에서 C3가 C2를 압도적으로 상회한다.}

\chg{다음으로 L2 회귀 검증(\jrev{H4}--\jrev{H5})을 보면,}
C1에서 $R^2 = 1.000$, C2에서 $R^2 = 0.998$, C3 전체에서 $R^2 = 0.858$이 달성되었다.
C3 전체의 $R^2$가 $0.95$ 미만이며, P85.8 분할 시
Core($R^2 = 0.999$)와 Tail($R^2 = 0.710$)의 결과를 보이므로,
\jrev{H4}는 Core 영역 기준에서 충족된다.
교호작용항의 유의성(\jrev{H5})은 $k_O \cdot O$가 C1의 $R^2$를
$0.869 \to 1.000$으로 향상시키는 결정적 기여를 통해 확인되었다.

\chg{이어 L3 강건성 검증(\jrev{H6}--\jrev{H7})을 보면,}
비중첩 대역 시드 체계의 독립성(\jrev{H6}, $|r|_{\max} = 0.016$)과
$10^8$회 다중 시드 시뮬레이션의 안정성(\jrev{H7}, 기각 0개)이 모두 확인되어,
seed $= 42$의 결과가 특정 시드에 의한 우연이 아님이 입증되었다.

\chg{끝으로 L4 이중채널 검증(\jrev{H8})을 보면,}
C3 Tail 영역의 $k_O$ 계수($-0.415$)가 Core 영역($-0.798$) 대비
$0.52$배로 감쇠되어, 극단 불확실성 영역에서 관측성의 복구 단축 효과가
약화됨이 확인되었다. 이 이중채널 감쇠 효과는 극단적 불확실성 하에서
불확실성 채널이 지배적으로 작용하여 관측성의 영향력이
상대적으로 축소됨을 의미하며, \jrev{\chref{ch:discussion}}에서 그 실무적 함의를
상세히 논의한다.


\bigskip
\jrev{이상의 시뮬레이션 결과와 가설 검증은 4.12절(교차 일관성)의 교차 조건 분석\chg{(}효과 크기 비교,
순서 가설 검증, 이중채널 감쇠 메커니즘 규명\chg{)}의 기초 자료로 활용된다.}


%% -------------------------------------------------------------------
\subsection{\jrev{\frev{분포 환경 적합성} — 종합 명제}}
\label{subsec:distribution_adequacy}
%% -------------------------------------------------------------------

\jrev{본 연구는 \citet{lee2026drto}에서 격자 탐색으로 \chg{도출하여 출판한} 곡률 매개변수 $\kappa^{*} = 1.2$를 새로운 분포 환경(\chg{C2 가우시안 이중채널과 C3 파레토 중꼬리})에 그대로 적용했다. 신규 격자 탐색을 통한 직접적 최적성 검증 대신, 본 절은 다음 9개 가설(H2--H10)의 PASS 결과를 통합 해석함으로써 $\kappa = 1.2$의 \frev{분포 환경 적합성(distribution-environment adequacy)을 시사한다}:}

\jrev{\chg{첫째, 물리적 안전성으로서 H2가 PASS하여 $0 < \DRTO < R_{\max} = 24$h가 C0--C3 모두에서 100\%
준수되며, $\kappa = 1.2$에서 분포 환경과 무관하게 비현실적 발산이 발생하지 않는다. 둘째, 이론적
일치로서 H3가 PASS하여 조건 간 순서 보존(C1$\leq$C0, C1$\leq$C2, P85.8 이상 C2$\leq$C3)이 이론 예측과
정합하며, $\kappa = 1.2$에서 분포 환경의 직관적 단조성이 유지된다. 셋째, 통계적 설명력으로서 H4가
PASS하여 C1 $R^2 = 1.000$, C2 $R^2 = 0.998$, C3 Core $R^2 = 0.999$로 $\kappa = 1.2$ 모델이 분포 환경
데이터를 잘 설명하고, H5가 PASS하여 $k_O \cdot O$의 $\Delta R^2 = 0.131$로 교호작용이 실제 발현한다.
넷째, 결과 안정성으로서 H6이 PASS하여 $|r|_{\max} = 0.016 < 0.02$로 비중첩 시드의 통계적 독립이
보장되고, H7이 PASS하여 $10^8$회 시뮬레이션에서 45개 회귀 계수가 모두 안정하여 $\kappa = 1.2$ 결과가
시드 우연의 산물이 아님이 확인된다. 다섯째, 메커니즘 발현으로서 H8이 PASS하여 C3 분포 환경에서
Tail/Core $k_O$ 비율이 $0.52 < 1$로 이중채널 비선형 감쇠 메커니즘이 $\kappa = 1.2$ 하에서 성공적으로
포착된다. 여섯째, 외부 검증으로서 H9와 H10이 PASS하여 정량 설문 평균 $\geq 4.0$, 정성 합의율
$\geq 0.80$으로 $\kappa = 1.2$ 모델 전체에 대한 BCM/DR 전문가 외부 검증을 통과한다(4.11절 참조).}}

\jrev{\textit{\citet{lee2026drto}의 $\kappa^{*} = 1.2$를 본 연구의 새로운 분포 환경(\chg{C2 가우시안과 C3 파레토})에 적용한 결과, 9개 가설(H2--H10)의 PASS가 \chg{i)} 물리적 안전성, \chg{ii)} 이론적 일치, \chg{iii)} 통계적 설명력, \chg{iv)} 결과 안정성, \chg{v)} 메커니즘 발현, \chg{vi)} 외부 검증의 \frev{6가지 측면에서 집합적으로 시사한다}. 따라서 $\kappa = 1.2$는 본 연구의 분포 환경에서 \frev{운영 적합(operationally adequate)하다}.}}

\jrev{\frev{분포 환경 적합성}\chg{의 학술적 한계로서,} \chg{\citet{lee2026drto}의} 4기준 종합점수 $F(\kappa) = f_1 \cdot f_2 \cdot f_3 \cdot f_4$ (\chg{실무 유의성, 비포화, 효과 크기, 설명력})는 C1 조건 기반으로 정의되었다. 따라서 \chg{C2와 C3} 조건에서 $F(\kappa)$를 직접 적용하여 새로운 격자 탐색을 수행하기에는 부적합하며, 이는 본 연구가 명시하는 학술적 한계이다. 그러나 위 9개 가설의 PASS는 격자 탐색의 직접적 대안을 제공하며, $\kappa = 1.2$의 \frev{분포 환경 적합성을 다중 차원에서 시사한다}.}

\fpara{``$\kappa^{*} = 1.2$가 본 연구의 분포 환경에서 운영 적합''이라는 \chg{종합 명제의 학술적 위상은 운영 적합성과 그 범위에 관한 것으로,} 위 ``종합 명제''에서 H2--H10 PASS의 통합 해석이 입증하는 바는 ``$\kappa = 1.2$ 사용 모델이 C2/C3 조건의 평가 기준(\chg{물리 경계, 이론 일치, 통계 설명력, 결과 안정성, 메커니즘 발현, 외부 검증})을 만족함''이며, 이는 \emph{운영 적합성}(operational adequacy)의 입증에 해당한다. 환경별 \emph{최적성}(optimality)의 직접 입증\chg{,} 즉 C2/C3 각각에서의 $F(\kappa)$ 격자 탐색을 통한 환경 무관 최적 $\kappa$ 도출은 본 연구의 연구 범위 밖이며 후속 연구의 과제로 남긴다. 이에 따라 본 연구는 $\kappa = 1.2$를 BCM 일반 운영의 \emph{\chg{잠정적이고 보수적인} 실무 기본값}으로 채택하며 (4.13.3절 참조), 본 채택은 운영 적합성 입증과 후속 절의 다중 보강 근거에 의해 뒷받침된다.}


%% -------------------------------------------------------------------
\fpara{\subsection{$\kappa = 1.2$의 잠정적 및 보수적 채택을 위한 정당화}}
\fpara{\label{subsec:kappa_provisional_justification}}
%% -------------------------------------------------------------------

\fpara{본 절은 \chg{\citet{lee2026drto}의} 도출값 $\kappa = 1.2$를 BCM 일반 운영의 \emph{\chg{잠정적이고 보수적인} 실무 기본값}(provisional and conservative operational default)으로 채택하는 정당성을 정리한다. 위 4.13절의 H2--H10 PASS는 \emph{운영 적합성}을 시사할 뿐 환경별 최적성을 입증하지 않으므로, 본 절은 \emph{\chg{잠정적이고 보수적인 기본값 채택}}을 뒷받침하는 보조 근거를 제시한다.}

\fpara{\chg{연구 범위와 적용 근거의 측면에서,} 본 연구는 C2/C3 조건에서 $F(\kappa)$의 직접 격자 탐색을 \emph{연구 범위 외}로 두고 C1 조건 도출 $\kappa = 1.2$를 BCM 일반 운영의 기본값으로 채택한다. 이 채택은 다음 근거에 의해 보강된다\chg{. 첫째,} $\kappa$\chg{는} sigmoid 입력 스케일링의 형태 모수로서 \chg{분포 평균과 분산}에 직접 의존하지 않으며, $z_U$의 무차원 정규화 인자 $R_0/(R_{\max} - R_0)$가 두 채널 effective span을 \chg{균형화한다. 둘째,} $\kappa = 1.2$\chg{는} 운영 신호 표본에서 $|z|$를 sigmoid \chg{비포화 비선형} 활성 영역($|z| \in [0.5, 3]$)에 \chg{집중시키며, 이는 입력 분포 모양과 독립이다. 셋째,} 4.11절 전문가 패널(Likert 평균 $\geq 4.0$, 합의도 $\geq 0.80$)이 도메인 신뢰성을 \chg{보강한다. 넷째,} ISO 22301:2019의 단일 정적값 결정 정신과 \chg{부합한다.}}

\fpara{\chg{잠정적이고 보수적인 사용의 실무적 필연성의 측면에서,} BCM 운영 실무는 어떤 $\kappa$ 값을 반드시 사용해야 하므로, 환경 무관 최적 $\kappa$의 존재 여부가 후속 연구에서 결정되기 전까지 \chg{\citet{lee2026drto}에서 가장 정밀하게 도출된 후보값} $\kappa = 1.2$를 잠정 채택하는 것이 합리적이다. 다른 값(예: 1.0, 1.5)은 도출 근거가 약하고, 모델 미사용은 실무 비현실이며, 무작위 선택은 정당화가 불가능하다. 따라서 본 채택은 \emph{학술적 최적성 입증}이 아니라 \emph{더 나은 대안이 부재한 상태에서의 최선 보수 선택}이라는 실무적 필연에 근거한다. 후속 연구에서 환경 무관 $\kappa$가 발견되면 갱신하고, 부재가 진단되면 환경별 적응 정책이 정립되며, 어느 경우든 본 잠정 위상은 보수적 운영 기준으로 유지된다.}

